\chapter{Conclusion}
\label{ch:conclusion}

This thesis added a biological constraint that \dg{} lacked, the foveated eye, and measured
its effect.
The arms were trained on MIT1003 against an otherwise-identical sharp-input control. At human acuity
the foveated input differs from sharp by at most $0.0071$~bits per fixation, under $1\%$ of what the
sharp control gains over the centerbias. A third set of arms, holding the
same blur fixed at the image centre, isolates gaze-contingency. The fixed-centre arms sit below the
gaze-contingent ones at every cutoff, by $0.003 \pm 0.003$~bits at human acuity and $0.013 \pm 0.004$ at
the coarsest, so the gap widens as the periphery is degraded further.

At human acuity the foveation removed nothing the fit depended on measurably. Past it, at the cutoffs
where the blur is stronger than a human eye applies, it did, within what a frozen backbone lets the
comparison attribute to the input (\secref{disc-limits}), and the cost appears as a loss of concentration
rather than of ranking.

Across images, \dg{}'s fit varies more than any arm differs from another, and that variation follows how
far the fifteen subjects spread; the arm comparison is small against it and is read within images for that
reason. The result rests on one dataset, one acuity profile and a backbone fitted to sharp photographs.

The foveation experiment leaves a working front-end.
The immediate open direction is replication, running the same comparison on a second adult dataset to
show whether the dose--response reported here is a property of the mechanism or of MIT1003. Beyond it
lies the developmental question, whether training a foveated model recapitulates, along its learning
trajectory, the child-to-adult changes in human gaze. This front-end was built to make that question
answerable.
